The Facility Location Problem (FLP) is a mathematical optimisation problem that aims 
to find the locations  a set of supply facilities can serve a set of demand sites at the lowest cost. 
Typically two types of cost are associated with the problem: 
\textit{setup costs} which are required to setup a facility and 
\textit{supply/connection costs} associated
with a facility meeting the demand of a given demand centre. Specific instances of the FLP 
are the Uncapacitated Facility Location Problem (UFLP) in which facilities have an unlimited supply of product, and 
the Capacitated Facility Location Problem (CFLP) in which facilities have a capacity on how much 
they can supply \cite{wu2006capacitated}. \par

Lots of challenges encountered in logistics and operations research lend themselves to 
formulation as CFLP problems. Canonical examples include determing locations for factories or warehouses 
to be established. Another set of problems that lend themselves to this formulation are those encountered in 
Energy Systems Modelling (ESM), a field employing the use of computational models to support policy
decisions regarding energy infrastructure. Such models typically aim at obtaining cost/benefit estimates
regarding changes to existing energy infrastructure or policy; for example, the consequences of subsidies, impact of new supply lanes 
and effects of introducing new energy generation and storage technologies. As many governments have
launched projects aiming to transition away from fossil fuels and replace these with renewables, 
energy systems modelling is now also employed in cost estimation, technology selection and pathway
estimation \cite{hoffmann2024review} \cite{nakata2011application}.\par

In this report we present a modification of the CFLP that is applied to determine optimal site choice 
for renewable energy infrastructure, which in this report we call the Energy Systems Location Problem (ESLP).
We then derive a Lagrangian Relaxation Heuristic (LR Heuristic) that is used to approximate solutions to this problem,
partially inspired by \cite{wu2006capacitated}, and apply it to a specific case: selecting optimal technologies and sites for UK 
renewable energy technologies that can meet the UK's energy demands. This report will be organised as follows. 
In Section \ref{sec:2} we first outline the CFLP, and then present a modified version to fit requirements for energy systems modelling 
and obtain the ESLP. We then derive the LR heuristic that will be tested. In Section \ref{sec:3} we describe the
case study we will test the ESLP and LR Heuristic on and how the data was obtained and processed.
We then present results, comparing a direct implementation of the ESLP to the derived LR Heuristic.


